\documentclass[11pt]{article}
\usepackage[margin=1.1in]{geometry}
\usepackage{amsmath,amssymb,amsthm}
\usepackage{hyperref}

\newtheorem{theorem}{Theorem}[section]
\newtheorem{lemma}[theorem]{Lemma}
\newtheorem{corollary}[theorem]{Corollary}
\theoremstyle{definition}
\newtheorem{definition}[theorem]{Definition}
\newtheorem{conjecture}[theorem]{Conjecture}
\theoremstyle{remark}
\newtheorem*{remark}{Remark}

\title{A Machine-Verified Disproof of TLMC Conjecture \#1243:\\ Rule 30's Central Column Is Not Square-Free}
\author{Batch 4 solution submission}
\date{September 12, 2026}

\begin{document}
\maketitle

\begin{abstract}
We disprove a conjecture of the Last Math Competition about the central column of Wolfram's rule~30 cellular automaton. The conjecture asserts that the column is ``square-free over all prefixes of length $2^k$'' and that ``among its factors of length $< 2^k$ none contains a square''. In the standard sense of combinatorics on words (a \emph{square} being any subword of the form $uu$ with $u$ nonempty), the column begins $1, 1, 0, 1, \ldots$, so its prefix of length $2^2 = 4$ already contains the square $11$. The disproof is by kernel-verified computation in Lean~4: the evaluation of the defining recursion of rule~30 is performed by \texttt{decide}, and the computation theorems depend on \emph{no} axioms at all.
\end{abstract}

\section{The conjecture}

\begin{conjecture}[TLMC-00000001243]\label{conj:1243}
The central column of Wolfram's rule~30 (the output sequence of column~0 from a single~1 as initial condition) is not eventually periodic, and among its factors of length $< 2^k$ none contains a square: the central column is square-free over all prefixes of length $2^k$.
\end{conjecture}

\section{The disproof}

\begin{definition}[Rule 30 on $\mathbb{Z}$]
Rows are functions $\mathbb{Z} \to \{0,1\}$; row $0$ is the indicator of $\{0\}$, and
\[
\operatorname{row}_{t+1}(p) \;=\; \operatorname{row}_t(p-1) \;\oplus\; \bigl(\operatorname{row}_t(p) \;\vee\; \operatorname{row}_t(p+1)\bigr),
\]
the local rule of automaton~30. The \emph{central column} is $c_t = \operatorname{row}_t(0)$.
\end{definition}

\begin{theorem}\label{thm:values}
The central column begins
\[
c_0\, c_1\, c_2\, c_3\, c_4\, c_5\, c_6\, c_7 \;=\; 1\, 1\, 0\, 1\, 1\, 1\, 0\, 0 .
\]
\end{theorem}

\begin{proof}
Direct evaluation of the defining recursion (Lean: \texttt{colPrefix\_8}, proved by \texttt{decide}, i.e.\ by the Lean kernel itself; the theorem depends on no axioms).
\end{proof}

\begin{theorem}\label{thm:1243}
Conjecture~\ref{conj:1243} is \textbf{false}. The prefix of length $2^2 = 4$ is $1\,1\,0\,1$, which contains the factor $11 = 1 \cdot 1$, a square (of period one). Hence the column is not square-free over the prefixes of length $2^k$; equally, a factor of length $< 2^2$ (namely $11$, of length $2$) is itself a square.
\end{theorem}

\begin{proof}
The factor $11$ occurs at positions $0$--$1$ of the prefix $1101$ from Theorem~\ref{thm:values}. A word of the form $uu$ with $u$ nonempty is a square; here $u = 1$. (Lean: \texttt{prefix4\_has\_square} and \texttt{tlm1243\_false}.)
\end{proof}

\begin{remark}
The non-eventual-periodicity part of Conjecture~\ref{conj:1243} is a famous open problem and is \emph{not} disputed here; the disproof concerns the square-freeness clause as literally stated. If the author intended ``no large squares'', the statement would need a lower bound on the square period, and as written the universal claim over all $k$ fails at $k = 2$.
\end{remark}

\section{Machine verification}

The formalization \texttt{Rule30.lean} (namespace \texttt{TLMCBatch4}) defines rule~30 on $\mathbb{Z}$ exactly as above, computes the first values by \texttt{decide}, and derives \texttt{tlm1243\_false}. Compilation with \texttt{lake env lean Rule30.lean} exits with code $0$ with zero errors and warnings; \texttt{\#print axioms} reports that \texttt{colPrefix\_8} and \texttt{colPrefix\_4} depend on \emph{no} axioms, and \texttt{tlm1243\_false} only on \texttt{[propext, Classical.choice, Quot.sound]}. No \texttt{sorry} appears.

\end{document}
